\documentclass[aspectratio=169,11pt]{beamer}

\usepackage[utf8]{inputenc}
\usepackage[T1]{fontenc}
\usepackage[spanish,es-nodecimaldot]{babel}
\usepackage{booktabs}
\usepackage{array}
\usepackage{graphicx}
\usepackage{listings}
\usepackage{xcolor}
\usepackage{tikz}
\usetikzlibrary{positioning}
\usepackage{ragged2e}
\usepackage{hyperref}

\definecolor{coursegreen}{HTML}{1B5E20}
\definecolor{coursegreenlight}{HTML}{E8F5E9}
\definecolor{coursegreenmid}{HTML}{43A047}
\definecolor{coursegray}{HTML}{455A64}
\definecolor{coursewhite}{HTML}{FFFFFF}

\usetheme{default}
\setbeamertemplate{navigation symbols}{}
\setbeamertemplate{frametitle}{
  \vspace{0.25cm}
  \insertframetitle\par
  \vspace{0.08cm}
  \color{coursegreenmid}\rule{\textwidth}{1.2pt}
}
\setbeamercolor{normal text}{fg=coursegray,bg=coursewhite}
\setbeamercolor{structure}{fg=coursegreen}
\setbeamercolor{frametitle}{fg=coursegreen}
\setbeamercolor{title}{fg=coursegreen}
\setbeamercolor{author}{fg=coursegray}
\setbeamercolor{block title}{fg=coursewhite,bg=coursegreen}
\setbeamercolor{block body}{fg=coursegray,bg=coursegreenlight}
\setbeamercolor{alertblock title}{fg=coursewhite,bg=coursegreenmid}
\setbeamercolor{alertblock body}{fg=coursegray,bg=coursegreenlight}
\setbeamercolor{exampleblock title}{fg=coursewhite,bg=coursegreen}
\setbeamercolor{exampleblock body}{fg=coursegray,bg=coursegreenlight}
\setbeamercolor{footline}{fg=coursegray,bg=coursewhite}

\setbeamertemplate{footline}{
  \leavevmode
  \hbox{
    \begin{beamercolorbox}[wd=.82\paperwidth,ht=2.5ex,dp=1ex,leftskip=1em]{footline}
      CourseSwap
    \end{beamercolorbox}
    \begin{beamercolorbox}[wd=.18\paperwidth,ht=2.5ex,dp=1ex,center]{footline}
      \insertframenumber/\inserttotalframenumber
    \end{beamercolorbox}
  }
}

\lstset{
  basicstyle=\ttfamily\scriptsize,
  keywordstyle=\color{coursegreen}\bfseries,
  commentstyle=\color{coursegray}\itshape,
  stringstyle=\color{coursegreenmid},
  frame=single,
  rulecolor=\color{coursegreenmid},
  backgroundcolor=\color{coursegreenlight},
  showstringspaces=false,
  breaklines=true,
  columns=fullflexible
}

\lstdefinelanguage{JavaScript}{
  morekeywords={const,let,var,function,return,if,else,await,async,new,import,from,
    export,default,throw,try,catch},
  sensitive=true,
  morecomment=[l]{//},
  morestring=[b]',
  morestring=[b]"
}

\title{CourseSwap}
\subtitle{De la necesidad académica al intercambio de grupos}
\author{\texorpdfstring{Entornos de Programación F1\\Alejandro Castro Rondon \and Jonk Keyler Sánchez Pabón \and Juan Carlos Elizalde Padilla}{Entornos de Programación F1 -- CourseSwap}}
\institute{Universidad Industrial de Santander}
\date{2026}

\begin{document}

\begin{frame}
  \titlepage
\end{frame}

\begin{frame}{¿Qué problema resuelve CourseSwap?}
  \begin{block}{Motivación}
    Cambiarse de grupo puede ser necesario cuando el horario deja de ser conveniente,
    pero encontrar a otra persona que quiera realizar el intercambio inverso exige
    búsquedas manuales y coordinación externa.
  \end{block}
  \begin{columns}[T]
    \column{0.48\textwidth}
      \textbf{Antes}
      \begin{itemize}
        \item Búsqueda informal de estudiantes.
        \item Validación manual de grupos y horarios.
        \item Coordinación dispersa y poco trazable.
      \end{itemize}
    \column{0.48\textwidth}
      \textbf{Con CourseSwap}
      \begin{itemize}
        \item Solicitudes con grupo actual y grupo deseado.
        \item Matching inverso por materia.
        \item Confirmación bilateral antes de ejecutar el cambio.
      \end{itemize}
  \end{columns}
\end{frame}

\begin{frame}{Flujo principal del producto}
  \begin{enumerate}
    \item El estudiante consulta sus materias, matrículas, grupos y horarios.
    \item Publica una solicitud de cambio.
    \item El sistema busca una solicitud inversa compatible.
    \item Se crea un match y se notifica a ambas personas.
    \item Cada participante confirma el intercambio.
    \item Con las dos confirmaciones, se intercambian las matrículas y se registra
          el cambio ejecutado.
  \end{enumerate}
  \vspace{0.25cm}
  \begin{alertblock}{Regla de seguridad funcional}
    El intercambio no se ejecuta con una sola confirmación: ambas partes deben
    aceptar el match.
  \end{alertblock}
\end{frame}

\begin{frame}{Funcionalidades implementadas}
\scriptsize
\begin{tabular}{>{\raggedright\arraybackslash}p{0.34\textwidth}
                >{\raggedright\arraybackslash}p{0.16\textwidth}
                >{\raggedright\arraybackslash}p{0.42\textwidth}}
\toprule
\textbf{Funcionalidad} & \textbf{Estado} & \textbf{Descripción} \\
\midrule
Registro de estudiantes & Implementado & Nombre, apellido, código, correo y contraseña con validación y BCrypt. \\
Inicio de sesión & Implementado & Spring Security, JWT y rutas protegidas. \\
Perfil autenticado & Implementado & Consulta del estudiante actual mediante \texttt{/estudiantes/me}. \\
Materias, grupos y horarios & Implementado & Consulta de catálogo y horarios para elegir destinos. \\
Consulta y creación de matrículas & Implementado & Gestión de la matrícula desde la interfaz. \\
Solicitudes de cambio & Implementado & Crear, editar, listar y cancelar solicitudes. \\
Conflictos de horario & Implementado & Validación al solicitar y al aceptar coincidencias. \\
Matching automático y manual & Implementado & Coincidencia inversa y búsqueda de solicitudes compatibles. \\
Confirmación bilateral & Implementado & Ambos estudiantes deben confirmar el match. \\
Ejecución e historial & Implementado & Intercambio de grupos y registro en \texttt{CambioEjecutado}. \\
Notificaciones por correo & Implementado & Brevo para solicitud creada, match encontrado e intercambio confirmado. \\
Tema y sesión & Implementado & Tema claro/oscuro, JWT persistido y guard de sesión. \\
\bottomrule
\end{tabular}
\end{frame}

\begin{frame}{Arquitectura general}
  \begin{center}
    \begin{tikzpicture}[node distance=0.8cm, every node/.style={font=\small},
      box/.style={draw=coursegreen, rounded corners, fill=coursegreenlight,
                  minimum width=3.1cm, minimum height=0.9cm, align=center},
      arrow/.style={->, thick, color=coursegreen}]
      \node[box] (frontend) {Frontend SPA\\React + Vite};
      \node[box, right=of frontend] (api) {API REST\\Spring Boot};
      \node[box, right=of api] (db) {Persistencia\\PostgreSQL};
      \node[box, below=of api] (mail) {Servicio de correo\\Brevo API};
      \draw[arrow] (frontend) -- node[above]{HTTP + JWT} (api);
      \draw[arrow] (api) -- node[above]{JPA/Hibernate} (db);
      \draw[arrow] (api) -- node[right]{REST + HTML} (mail);
    \end{tikzpicture}
  \end{center}
  \vspace{0.25cm}
  \begin{block}{Backend}
    API REST monolítica por capas: Controller $\rightarrow$ Service $\rightarrow$
    Repository $\rightarrow$ PostgreSQL. DTOs y un manejador global separan
    contratos, persistencia y errores de negocio.
  \end{block}
\end{frame}

\begin{frame}{El nuevo frontend: una SPA modular}
  El proyecto conserva un frontend estático legado, pero la experiencia principal
  se reorganizó en una SPA React con rutas, hooks y componentes reutilizables.
  \vspace{0.25cm}
  \begin{columns}[T]
    \column{0.48\textwidth}
      \begin{block}{Características}
        \begin{itemize}
          \item Renderizado declarativo con React.
          \item Navegación sin recargar la página.
          \item Hooks para autenticación, solicitudes, matrículas y matches.
          \item Componentes MUI para una interfaz consistente.
          \item Estado remoto administrado por TanStack Query.
        \end{itemize}
      \end{block}
    \column{0.48\textwidth}
      \begin{block}{Separación de responsabilidades}
        \begin{itemize}
          \item \texttt{pages/}: composición de pantallas.
          \item \texttt{components/}: UI reutilizable.
          \item \texttt{hooks/}: lógica de datos y mutaciones.
          \item \texttt{services/}: llamadas HTTP.
          \item \texttt{store/}: notificaciones globales.
        \end{itemize}
      \end{block}
  \end{columns}
\end{frame}

\begin{frame}[fragile]{Árbol de carpetas del frontend}
\scriptsize
\begin{lstlisting}
frontend-spa/
+-- public/
+-- src/
|   +-- api/              # Axios y configuracion del JWT
|   +-- assets/           # Logo y recursos visuales
|   +-- components/       # Formularios, tarjetas, layout y tablas
|   +-- constants/        # Constantes de agenda y UI
|   +-- hooks/            # useAuth, useQuery y useMutation
|   +-- pages/            # Auth, Dashboard, Solicitudes y Matriculas
|   +-- services/         # auth, match, materias, matricula y solicitudes
|   +-- store/            # notifyStore con Zustand
|   +-- App.jsx           # Router y rutas protegidas
|   +-- main.jsx          # React root + QueryClientProvider
+-- index.html
+-- package.json
+-- vite.config.js
+-- eslint.config.js
\end{lstlisting}
\end{frame}

\begin{frame}{¿Para qué sirve Vite?}
  \begin{block}{Definición}
    Vite es la herramienta de desarrollo y build del frontend. Proporciona un
    servidor local rápido durante el desarrollo y genera los archivos optimizados
    para producción.
  \end{block}
  \begin{columns}[T]
    \column{0.48\textwidth}
      \textbf{Durante el desarrollo}
      \begin{itemize}
        \item Servidor con Hot Module Replacement (HMR).
        \item Recarga rápida de módulos modificados.
        \item Integración directa con React mediante \texttt{@vitejs/plugin-react}.
      \end{itemize}
    \column{0.48\textwidth}
      \textbf{Para producción}
      \begin{itemize}
        \item Resuelve imports ES Modules.
        \item Minifica y empaqueta los recursos.
        \item Genera la carpeta \texttt{dist/} para desplegar.
      \end{itemize}
  \end{columns}
\end{frame}

\begin{frame}{Dependencias principales del frontend}
\scriptsize
\begin{tabular}{>{\raggedright\arraybackslash}p{0.28\textwidth}
                >{\raggedright\arraybackslash}p{0.60\textwidth}}
\toprule
\textbf{Dependencia} & \textbf{Propósito en CourseSwap} \\
\midrule
\texttt{react}, \texttt{react-dom} & Construcción declarativa de la interfaz y montaje de la aplicación. \\
\texttt{vite} & Servidor de desarrollo, HMR y build de producción. \\
\texttt{react-router-dom} & Rutas, navegación client-side, layout compartido y rutas protegidas. \\
\texttt{@mui/material} & Componentes visuales accesibles: botones, formularios, tarjetas, diálogos y grids. \\
\texttt{@mui/icons-material} & Iconos consistentes para estados, acciones y navegación. \\
\texttt{@emotion/react}, \texttt{@emotion/styled} & Motor de estilos usado por Material UI. \\
\texttt{@tanstack/react-query} & Queries, mutaciones, caché, estados de carga e invalidación. \\
\texttt{axios} & Cliente HTTP común e interceptor para adjuntar el JWT. \\
\texttt{zustand} & Estado global pequeño para notificaciones. \\
\texttt{eslint} y plugins & Detección de errores y consistencia durante el desarrollo. \\
\bottomrule
\end{tabular}
\end{frame}

\begin{frame}{¿Por qué TanStack Query?}
  \begin{block}{Problema que resuelve}
    El frontend consulta datos remotos que tienen estados propios: cargando,
    éxito, error, datos antiguos y actualización después de una acción.
    TanStack Query encapsula ese ciclo sin convertir todo en estado manual.
  \end{block}
  \begin{itemize}
    \item Centraliza caché por claves como \texttt{['solicitudes', studentId]}.
    \item Evita solicitudes duplicadas y conserva datos recientes.
    \item Expone estados claros: \texttt{isPending}, \texttt{isError} y \texttt{data}.
    \item Permite invalidar y refrescar datos después de crear, editar o confirmar.
    \item Separa el estado remoto del estado visual local de los formularios.
  \end{itemize}
\end{frame}

\begin{frame}{Queries y mutaciones}
  \begin{columns}[T]
    \column{0.48\textwidth}
      \begin{block}{Query: lectura}
        Una query obtiene datos del servidor y los mantiene disponibles en caché.
        \[
          \texttt{useQuery} \Rightarrow
          \text{data + loading + error}
        \]
        \vspace{-0.1cm}
        Ejemplos:
        \begin{itemize}
          \item Solicitudes del estudiante.
          \item Matches activos.
          \item Horarios de una materia.
        \end{itemize}
      \end{block}
    \column{0.48\textwidth}
      \begin{block}{Mutación: escritura}
        Una mutación ejecuta un cambio en el servidor y luego actualiza las
        queries relacionadas.
        \[
          \texttt{useMutation} \Rightarrow
          \text{acción + éxito + error}
        \]
        \vspace{-0.1cm}
        Ejemplos:
        \begin{itemize}
          \item Crear o editar una solicitud.
          \item Confirmar un match.
          \item Eliminar una solicitud.
        \end{itemize}
      \end{block}
  \end{columns}
\end{frame}

\begin{frame}[fragile]{Ejemplo: mutación e invalidación}
\begin{lstlisting}[language=JavaScript]
const updateSolicitudMutation = useMutation({
  mutationFn: ({ solicitudId, nuevoGrupoId }) =>
    solicitudesService.updateSolicitud(solicitudId, nuevoGrupoId),
  onSuccess: () => {
    queryClient.invalidateQueries({
      queryKey: ['solicitudes', studentId],
    })
  },
})
\end{lstlisting}
\begin{exampleblock}{Idea clave}
  La interfaz no necesita reconstruir manualmente toda la lista: se invalida la
  query y TanStack Query vuelve a sincronizarla con el backend.
\end{exampleblock}
\end{frame}

\begin{frame}{Uso de react-router-dom}
  \texttt{react-router-dom} permite que la SPA cambie de pantalla sin recargar
  completamente el documento.
  \begin{itemize}
    \item \texttt{createBrowserRouter}: define el árbol de rutas.
    \item \texttt{Outlet}: renderiza las páginas dentro del layout principal.
    \item \texttt{NavLink}: enlaces con estado activo en la navbar.
    \item \texttt{Navigate}: redirección cuando no existe una sesión válida.
    \item \texttt{ProtectedRoute}: impide acceder a la aplicación sin autenticación.
  \end{itemize}
  \vspace{0.25cm}
  \begin{block}{Resultado}
    El layout, la navbar y la autenticación permanecen centralizados mientras
    cada página se concentra en su propia responsabilidad.
  \end{block}
\end{frame}

\begin{frame}{Material UI en la interfaz}
  \begin{block}{¿Qué es Material UI?}
    Material UI (MUI) es una biblioteca de componentes React basada en los
    principios de Material Design. Proporciona componentes listos para usar,
    personalizables y con patrones de interacción consistentes.
  \end{block}
  \begin{columns}[T]
    \column{0.48\textwidth}
      \textbf{Componentes utilizados}
      \begin{itemize}
        \item \texttt{Button}, \texttt{TextField}, \texttt{Select}.
        \item \texttt{Card}, \texttt{Paper}, \texttt{Dialog}.
        \item \texttt{Alert}, \texttt{Chip}, \texttt{Grid}, \texttt{Stack}.
        \item \texttt{AppBar}, \texttt{Drawer}, \texttt{Tabs}.
      \end{itemize}
    \column{0.48\textwidth}
      \textbf{Beneficios}
      \begin{itemize}
        \item Diseño coherente entre pantallas.
        \item Responsive design con breakpoints.
        \item Accesibilidad y estados visuales incorporados.
        \item Personalización mediante \texttt{sx} y la paleta del tema.
        \item Menos CSS repetido y mayor velocidad de desarrollo.
      \end{itemize}
  \end{columns}
\end{frame}

\begin{frame}{Iconos: significado visual y consistencia}
  \begin{block}{¿Por qué usar los iconos de MUI?}
    \texttt{@mui/icons-material} ofrece iconos alineados con Material Design y
    compatibles con los componentes de la biblioteca.
  \end{block}
  \begin{itemize}
    \item Estados de solicitudes: propuesta, match, confirmada y cancelada.
    \item Acciones: confirmar, cancelar, editar y crear.
    \item Navegación: menú móvil, inicio y accesos rápidos.
    \item Refuerzo visual de alertas y estados vacíos.
  \end{itemize}
  \vspace{0.2cm}
  \begin{exampleblock}{Beneficio UX}
    El icono no reemplaza el texto: lo refuerza. Esto mejora el escaneo visual y
    hace que las acciones importantes sean reconocibles de inmediato.
  \end{exampleblock}
\end{frame}

\begin{frame}{Servicio de correo con Brevo}
  \begin{block}{Objetivo}
    Informar al estudiante de eventos importantes sin obligarlo a permanecer
    dentro de la aplicación: creación de solicitud, match encontrado y
    confirmación completa del intercambio.
  \end{block}
  \begin{enumerate}
    \item El servicio de dominio detecta el evento.
    \item Construye un contexto con nombre, materia, grupos, código y match.
    \item Thymeleaf procesa una plantilla HTML.
    \item \texttt{EmailService} envía el HTML a la API SMTP de Brevo.
    \item Brevo entrega el correo al destinatario.
  \end{enumerate}
  \vspace{0.15cm}
  \texttt{POST https://api.brevo.com/v3/smtp/email}
\end{frame}

\begin{frame}[fragile]{Implementación del EmailService}
\begin{lstlisting}[language=Java]
@Async("emailExecutor")
public void sendEmail(String destinationEmail,
                      String subject,
                      String templateName,
                      Context context) {
    String htmlBody = templateEngine.process(templateName, context);
    Map<String, Object> request = Map.of(
        "sender", Map.of("email", senderEmail),
        "to", new Map[]{Map.of("email", destinationEmail)},
        "subject", subject,
        "htmlContent", htmlBody
    );
    restClient.post()
        .uri(BREVO_SEND_EMAIL_URL)
        .header("api-key", brevoApiKey)
        .contentType(MediaType.APPLICATION_JSON)
        .body(request)
        .retrieve()
        .toBodilessEntity();
}
\end{lstlisting}
\end{frame}

\begin{frame}{Decisiones técnicas del correo}
  \begin{columns}[T]
    \column{0.48\textwidth}
      \begin{block}{Asincronía}
        \texttt{@Async("emailExecutor")} desacopla el envío de la respuesta
        principal. La creación de una solicitud o la confirmación no queda
        bloqueada esperando a Brevo.
      \end{block}
      \begin{block}{Plantillas}
        Thymeleaf genera HTML con datos del evento y mantiene la presentación
        del correo fuera de la lógica del servicio.
      \end{block}
    \column{0.48\textwidth}
      \begin{block}{Configuración}
        Las credenciales se obtienen mediante variables de entorno:
        \texttt{BREVO\_API\_KEY} y \texttt{BREVO\_SENDER\_EMAIL}.
      \end{block}
      \begin{block}{Resiliencia}
        Se validan credenciales antes del envío y se registran errores de
        \texttt{RestClient} con el destinatario y la causa.
      \end{block}
  \end{columns}
\end{frame}

\begin{frame}{¿Cuándo se envían correos?}
\small
\resizebox{\textwidth}{!}{%
\begin{tabular}{>{\raggedright\arraybackslash}p{0.29\textwidth}
                >{\raggedright\arraybackslash}p{0.26\textwidth}
                >{\raggedright\arraybackslash}p{0.31\textwidth}}
\toprule
\textbf{Evento} & \textbf{Plantilla} & \textbf{Destinatario y propósito} \\
\midrule
Solicitud sin match & \texttt{solicitud-creada.html} & El estudiante recibe confirmación de creación, materia y grupos. \\
Match automático & \texttt{match-encontrado.html} & Ambos estudiantes reciben el código de la otra persona y deben revisar el match. \\
Match manual & \texttt{match-encontrado.html} & Se notifica a los dos participantes cuando se genera la coincidencia. \\
Confirmación bilateral & \texttt{match-confirmado.html} & Ambos reciben la confirmación del intercambio ejecutado. \\
\bottomrule
\end{tabular}
}
\vspace{0.25cm}
\begin{alertblock}{Importante}
  El correo acompaña el flujo de negocio; la fuente de verdad continúa siendo
  el estado persistido en la API y la base de datos.
\end{alertblock}
\end{frame}

\begin{frame}{Seguridad y despliegue}
  \begin{itemize}
    \item JWT para autenticación stateless y protección de endpoints.
    \item Axios agrega automáticamente el token en las solicitudes de la SPA.
    \item Variables de entorno para secretos de JWT, base de datos y Brevo.
    \item Backend desplegado en Render con PostgreSQL.
    \item Frontend SPA desplegado en GitHub Pages mediante GitHub Actions.
    \item Swagger/OpenAPI disponible para explorar el contrato de la API.
  \end{itemize}
  \vspace{0.25cm}
  \begin{block}{Principio aplicado}
    Los secretos no forman parte del código fuente ni de la configuración
    pública del frontend.
  \end{block}
\end{frame}

\begin{frame}{Cierre}
  \begin{center}
    {\Huge\color{coursegreen}\textbf{CourseSwap}}
    
    \vspace{0.4cm}
    {\large Una plataforma para convertir una búsqueda manual en un intercambio
    seguro, trazable y confirmado por ambas partes.}
    
    \vspace{0.7cm}
    \textcolor{coursegreenmid}{Gracias}
  \end{center}
\end{frame}

\end{document}
